\section{Introduction}

\section{Prior Work}

\subsection{Shor's Algorithm and Demonstrations}
\textcolor{red}{@Infleqtion: Have already started adding}

\subsection{Quantum Error Detection}
\textcolor{blue}{Optional, can add the background from the prior material sciences paper}

\section{Sqale Architecture}
\textcolor{red}{@Infleqtion: details about the architecture}

\subsection{Instruction Set Architecture}
\textcolor{red}{@Infleqtion: Describe the device level instructions, including movement capabilities. Also include corresponding fidelities motivating exploration of different iceberg mappings.}


\section{Iceberg Code}
The simplest family of error detecting codes is $[[n, n-2, 2]]$ family of codes, also known as the Iceberg code\cite{}. This family can detect up to one $X$ or $Z$ error on any of the $n$ physical qubits. The $X$ and $Z$ stabilizers are given by $X_1 \otimes X_2 \cdots X_n$ and $Z_1 \otiesm Z_2 \cdots Z_n$, which detects a $Z$ and $X$ error, respectively. In addition to this, the Iceberg family has a universal logical gate set\cite{}, making it a promising candidate for near-term error correction quantum devices. The universal gate set can also be made fault-tolerant using the techniques proposed in~\cite{}, which we discuss further in Section~\ref{sec:}. \par
Although the encoding rate of the code ($= (n-2)/n$) approaches $1$ for large $n$, the errors detected per physical qubit approach $0$. Therefore, the more interesting codes from this family are the smaller codes, such as the $[[4, 2, 2]]$ code. The $[[4, 2, 2]]$ code, for example, can detect up to $1$ error for $2$ physical qubits, and only uses $1$ extra physical qubit for each logical qubit.

\subsection{(Fault-Tolerant) Universal Gate Set}
\textcolor{brown}{@Baker group: list out the universal gate set and their fault-tolerant implementation.}

\subsection{Mapping to the Iceberg Code}
\textcolor{brown}{@Baker group: give motivation using some example circuit (probably a 2-qubit circuit): show unencoded vs encoded. For encoded, show different possible mappings and how they significantly vary the TVD.}


\section{ACRONYM Toolkit}
\subsection{Notation}
\textcolor{brown}{@Baker group: talk about the distinction between program circuit, logical circuit, physical (program) circuit and atom level.}

\subsection{Program to Physical Compilation}
\subsubsection{Unique Permutations}
\textcolor{brown}{@Baker group: define what are unique permutations and how to generate them}

\subsubsection{Translation}
\textcolor{brown}{@Baker group: show the pipeline of the toolkit which choses each unique permutations and translates from the program circuit to the logical circuit. Specifically, talk about the choices made during translation.}

\subsubsection{Logical to Physical Translation}
We build a circuit by replacing each logical gate with its physical circuit implementation (from Section~\ref{sec:}). Note that this is still agnostic to the placement on the atom grid and the movements needed to execute the gates.

\subsection{Physical to Atom Compilation}
\textcolor{red}{@Infleqtion: details about the atom placement and movement strategies.}
\textcolor{brown}{@Baker group: talk about challenges with automating this and how we solve them(?)}

\subsection{Noisy Simulation and Decision}
We use sqalesim\cite{} to simulate the different circuits generated for each unique permutation. We also run a noiseless unencoded simulation to generate the oracle probability distribution. We then compute the Total Variational Distance (TVD) between each permutation circuit and the oracle. Finally, we choose the circuit which gives the smallest TVD\@.
